%% “华为杯”第二十三届中国研究生数学建模竞赛 论文 —— 主文件
%% 编译：latexmk -xelatex main.tex   （或 xelatex main.tex 两遍）
%%
%% 类选项：coverpagenumber（封面显示页码 0）、nocover（无封面）
%%
%% ⚠ 隐私提示：下面的 \school / \teamnumber / \memberA..C 默认是占位符。
%%   · 若本仓库公开，提交前请务必保持占位符，不要写入真实学校、队号、姓名；
%%   · 真实信息只在本机填写并本地编译，PDF 成品单独提交，不要推送到公开仓库。
\documentclass{gmcm}

% ------------------------- 封面 / 摘要页信息 -------------------------
\school{你的大学名称}                 % 学校
\teamnumber{参赛队号}                 % 参赛队号（提交前填写）
\memberA{队员1}                      % 队员 1（提交前填写）
\memberB{队员2}                      % 队员 2（提交前填写）
\memberC{队员3}                      % 队员 3（提交前填写）
\title{研究生 数学建模 题目} % 题目（写在摘要页“题目：”之后）
\keywords{关键词1；关键词2；关键词3；关键词4}

% ------------------------- 自定义宏 -------------------------
\newcommand{\E}{\mathbb{E}}
\newcommand{\dd}{\mathrm{d}}

\begin{document}

\makecover                          % 封面（不计页码）

\begin{abstract}                    % 摘要页：页码从 1 开始
% 摘要：简明扼要，需包含 建模思路、主要方法、模型、结果与结论、创新点
% 篇幅一般不超过两页，无需译成英文。评阅时摘要与论文都会审阅。

本文针对……问题，建立了……模型。首先，……（建模思路）。

针对问题一，……（主要方法与模型）。求解得到……（结果与结论）。

针对问题二，……

本文的创新点在于：（1）……；（2）……；（3）……。

\end{abstract}

\tableofcontents                    % 目录：摘要、目录、各级标题、参考文献、附录

% ------------------------- 正文（一级标题 四号黑体居中） -------------------------
\section{问题重述}

\subsection{问题背景}
随着人工智能技术的快速发展，具身智能、智能交互和人机协同等应用逐渐从单纯的信息识别向更加深入的人类状态理解发展。在实际人机交互过程中，人的情感状态往往直接影响其语言表达、语音语调、面部表情以及肢体行为。因此，对人类情感进行准确识别与预测，是提高智能系统交互自然性和理解能力的重要基础。

与传统的单模态情感分析不同，真实环境中的情感表达通常由多种信息共同构成。文本模态主要反映语言内容及其语义信息，语音模态能够体现语调、节奏、音量等情感变化，视觉模态则可以通过面部表情和行为动作反映个体的外在情绪状态。不同模态之间具有较强的互补性，综合利用文本、语音和视觉信息，有助于更加全面地刻画人的情感状态。

然而，复杂场景下的多模态情感预测仍面临较大困难。一方面，文本、语音和视觉数据具有不同的采样频率、特征形式和时间尺度，不同模态的信息难以直接建立统一的时序对应关系；另一方面，在实际数据采集过程中，可能受到环境噪声、画面遮挡和语音转写异常等因素影响，使部分时段的模态信息缺失，从而降低传统多模态融合模型的稳定性。此外，现有情感预测模型通常更加关注最终预测结果，而对不同模态在决策过程中的作用以及关键情感证据缺乏直观、可量化的解释。

\subsection{问题重述}

\begin{enumerate}[label=\textbf{问题\arabic*:},leftmargin=*]
    \item 多模态特征提取与时序对齐,
从原始视频中分别提取文本、语音和视觉三模态情感特征，并建立统一的时序对应关系。保证各模态特征与原始样本、时间位置一一对应，为后续情感预测提供规范、可复现的输入。
    \item 模态局部缺失下的鲁棒情感预测,
针对文本、语音或视觉在部分连续时间段缺失的情况，建立能够动态利用剩余有效模态的鲁棒预测模型。同时分析缺失模态类型、缺失位置和缺失时长对情感极性分类与情感强度回归性能的影响。
    \item 可解释性情感预测,
建立既能输出情感极性和情感强度，又能解释预测依据的多模态模型。量化文本、语音、视觉三种模态的贡献，并定位对预测起关键作用的文本片段、语音时段和视觉关键帧。
\end{enumerate}
     % 问题重述
\section{问题分析}

\subsection{问题一的分析}

……

\subsection{问题二的分析}

……
\subsection{问题三的分析}

……        % 问题分析
\section{模型假设与符号说明}

\subsection{模型假设}

\begin{enumerate}[label=假设\arabic*：,leftmargin=*]
  \item ……
  \item ……
\end{enumerate}

\subsection{符号说明}

\begin{table}[H]
  \centering
  \caption{主要符号说明}
  \begin{tabular}{clc}
    \toprule
    符号 & 含义 & 单位 \\
    \midrule
    $x_i$    & …… & -- \\
    $\theta$ & …… & rad \\
    \bottomrule
  \end{tabular}
\end{table}
     % 模型假设与符号说明
\section{问题一模型的建立与求解}

\subsection{模型建立}

目标函数为
\begin{equation}
  \min_{x}\; f(x)=\sum_{i=1}^{n} c_i x_i ,
  \qquad \text{s.t.}\quad Ax\le b,\; x\ge 0 .
  \label{eq:q1-obj}
\end{equation}
式~\eqref{eq:q1-obj} 中……。相关方法参见文献\cite{ref-book}\cite{ref-journal}。

\subsection{模型求解}

% 图示例（图片放在 figures/ 目录）
% \begin{figure}[htbp]
%   \centering
%   \includegraphics[width=0.7\linewidth]{figures/example}
%   \caption{……}
%   \label{fig:q1-example}
% \end{figure}

\subsection{结果分析}

……

% 若数据分析结果参考了 AI 工具输出，须在其前导或后面添加注释（附件4 第4条）：
\ainote{本节数据分析结果的部分处理思路参考了 \aitool{XXX}{X.X}{XXX 公司}{2026-XX-XX} 的输出，经人工核验与复算。}
        % 问题一 模型建立与求解
\section{问题二模型的建立与求解}

\subsection{模型建立}

……

\subsection{模型求解}

……
\subsection{结果分析}

……
        % 问题二 模型建立与求解
\section{问题三模型的建立与求解}

\subsection{模型建立}

……

\subsection{模型求解}

……
\subsection{结果分析}

……
        % 问题三 模型建立与求解
\section{模型的评价与推广}

\subsection{模型的优点}

……

\subsection{模型的不足}

……

\subsection{模型的推广}

……
      % 模型评价与推广

% ------------------------- 参考文献（按格式规范手写） -------------------------
% 参考文献：按正文引用次序排列；正文引用处用方括号标示编号，如 \cite{ref-book}。
% 引用书籍必须指出页码。三类表述方式（附件2）：
%   书籍：[编号] 作者，书名，出版地：出版社，起止页码，出版年。
%   期刊：[编号] 作者，论文名，杂志名，卷期号：起止页码，出版年。
%   网络：[编号] 作者，资源标题，网址，访问时间（年月日）。
% AI 工具（附件4 第3条，仅用于无法确认来源的模型/公式，存在不被认可的风险）：
%   [编号] 人工智能工具名称，版本/型号，开发机构/公司，版本发布日期。
\begin{thebibliography}{99}
  \bibitem{ref-book}    姜启源，谢金星，叶俊，数学模型（第五版），北京：高等教育出版社，123--145，2018。
  \bibitem{ref-journal} 张三，李四，基于……的……模型，系统工程理论与实践，40(3)：567--578，2020。
  \bibitem{ref-web}     王五，……数据集说明，https://example.com/dataset，2026年9月16日。
  % \bibitem{ref-ai}    XXX（人工智能工具），版本/型号 X.X，XXX 公司，2026年X月X日。
\end{thebibliography}


% ------------------------- 附录 -------------------------
\clearpage
\appendix
\section{支撑材料清单}

\begin{enumerate}[leftmargin=*]
  \item \texttt{code/solve\_q1.py}：问题一求解程序；
  \item \texttt{code/solve\_q2.py}：问题二求解程序；
  \item \texttt{data/}：计算结果。
\end{enumerate}

\section{主要程序源代码}

% 若使用 AI 工具辅助编程，程序前面必须添加注释（附件4 第5条），见 code/solve_q1.py 开头。
\lstinputlisting[language=Python,caption={问题一求解程序 solve\_q1.py}]{code/solve_q1.py}


\end{document}
